% golden_baseline_settings.tex
%
% Companion listing to Appendix C of the thesis. The thesis prints only the
% provenance of the sigma-bearing settings; the full resolved listing lives here.
%
% The table body is held in _settings_table_body.tex so that a regenerated
% listing can replace it without touching this wrapper.
%
% Build: pdflatex golden_baseline_settings   (twice, for the longtable widths)

\documentclass[10pt,a4paper]{article}
\usepackage[T1]{fontenc}
\usepackage{lmodern}
\usepackage[margin=2cm]{geometry}
\usepackage{array}
\usepackage{longtable}
\usepackage{setspace}
\pagestyle{plain}

\begin{document}

\begin{center}
{\large\bfseries Golden-baseline reference scenario, resolved settings}\\[4pt]
Companion to Appendix C of the thesis.\\
Generated from \texttt{config/golden\_baseline.json}.
\end{center}

\vspace{1em}

\noindent
The settings below are the resolved contents of
\texttt{config/golden\_baseline.json}, excepting eight that name output files and set
diagnostic storage and so cannot reach any reported quantity. Two cautions carry with
the listing. First, the file is an overlay, so it states only what differs from the
master configuration and every setting not listed here takes its master default.
Second, a setting appearing here means the scenario file declares it, not that it
necessarily reaches the physics. The framework carries keys that no longer have a
consumer, and a listing cannot by itself distinguish a live setting from an inert one.
Claims in the thesis rest on measured runs and on the regression gates, never on the
presence of a key in a configuration file.

\vspace{1em}

\begin{center}
\small
\singlespacing
\begin{longtable}{>{\raggedright\arraybackslash\scriptsize\ttfamily}p{0.27\textwidth}|>{\raggedright\arraybackslash}p{0.38\textwidth}|>{\raggedright\arraybackslash\scriptsize\ttfamily}p{0.21\textwidth}}
\caption{Settings of the golden-baseline reference scenario}
\label{tab:golden_baseline_settings} \\
\hline
\textbf{\normalfont\normalsize Setting} & \textbf{\normalsize Short description} & \textbf{\normalfont\normalsize Value} \\ \hline
\endfirsthead
\multicolumn{3}{l}{\emph{\normalfont\normalsize Table \ref{tab:golden_baseline_settings}, continued from the previous page}} \\ \hline
\textbf{\normalfont\normalsize Setting} & \textbf{\normalsize Short description} & \textbf{\normalfont\normalsize Value} \\ \hline
\endhead
\hline \multicolumn{3}{r}{\emph{\normalfont\normalsize continued on the next page}} \\
\endfoot
\hline
\endlastfoot
\multicolumn{3}{l}{\textbf{Scenario and run}} \\ \hline
\texttt{scenario.\allowbreak n\allowbreak Receivers} & Number of receive antenna phase centres on the spacecraft & \texttt{4} \\
\texttt{scenario.\allowbreak n\allowbreak Space\allowbreak Assets} & Number of space assets simulated & \texttt{1} \\
\texttt{scenario.\allowbreak n\allowbreak Towers} & Number of ground tracking stations in the network & \texttt{5} \\
\texttt{scenario.\allowbreak name} & Scenario name label used in reports and run fingerprints & \texttt{golden\_\allowbreak baseline} \\
\texttt{scenario.\allowbreak orbit\allowbreak Class} & Orbit class of the tracked spacecraft & \texttt{GEO} \\
\texttt{simulation.\allowbreak dt\_\allowbreak s} & Measurement and propagation time step in seconds & \texttt{1} \\
\texttt{simulation.\allowbreak seed} & Master random number seed for bit-reproducible runs & \texttt{42} \\
\hline
\multicolumn{3}{l}{\textbf{Grade overlays}} \\ \hline
\texttt{atmosphere.\allowbreak estimate\allowbreak Iono} & Profile shortcut enabling slant ionosphere states, dropping Klobuchar correction & \texttt{false} \\
\texttt{atmosphere.\allowbreak ionosphere\allowbreak Free} & Selects ionosphere-free code combination rather than raw dual frequency & \texttt{false} \\
\texttt{atmosphere.\allowbreak realistic} & Selects the physically realistic troposphere and ionosphere profile & \texttt{true} \\
\texttt{atmosphere.\allowbreak shared\allowbreak Across\allowbreak Antennas.\allowbreak enable} & Shares one scintillation draw across all receive antennas & \texttt{true} \\
\texttt{realism.\allowbreak grade} & Opt-in switch for the whole realism error overlay & \texttt{true} \\
\texttt{realism.\allowbreak include.\allowbreak inter\allowbreak Antenna\allowbreak Carrier\allowbreak Bias} & Overlay sub-toggle for per-antenna carrier line bias & \texttt{true} \\
\texttt{realism.\allowbreak include.\allowbreak luni\allowbreak Solar} & Overlay sub-toggle for the reduced-dynamics luni-solar stress case & \texttt{false} \\
\texttt{realism.\allowbreak include.\allowbreak solid\allowbreak Earth\allowbreak Tide} & Overlay sub-toggle for truth-only solid Earth tide displacement & \texttt{false} \\
\hline
\multicolumn{3}{l}{\textbf{Orbit and force model}} \\ \hline
\texttt{orbit.\allowbreak truth.\allowbreak perturbations.\allowbreak epoch\allowbreak JD\_\allowbreak TT} & Sun and Moon ephemeris reference epoch, Julian date & \texttt{2451545.\allowbreak 0} \\
\texttt{orbit.\allowbreak truth.\allowbreak perturbations.\allowbreak luni\allowbreak Solar.\allowbreak enable} & Sun and Moon third-body acceleration in truth propagation & \texttt{true} \\
\texttt{orbit.\allowbreak truth.\allowbreak perturbations.\allowbreak srp.\allowbreak Cr} & Cannonball radiation pressure reflectivity coefficient, dimensionless & \texttt{1.\allowbreak 3} \\
\texttt{orbit.\allowbreak truth.\allowbreak perturbations.\allowbreak srp.\allowbreak area\allowbreak To\allowbreak Mass\_\allowbreak m2pkg} & Spacecraft area-to-mass ratio, square metres per kilogram & \texttt{0.\allowbreak 02} \\
\texttt{orbit.\allowbreak truth.\allowbreak perturbations.\allowbreak srp.\allowbreak enable} & Solar radiation pressure acceleration in truth propagation & \texttt{true} \\
\texttt{orbit.\allowbreak truth.\allowbreak perturbations.\allowbreak srp.\allowbreak shadow} & Eclipse shadow geometry used for radiation pressure & \texttt{cylindrical} \\
\hline
\multicolumn{3}{l}{\textbf{Propagation and geometry effects}} \\ \hline
\texttt{effects.\allowbreak antenna.\allowbreak pcv\allowbreak Model} & Antenna phase centre variation model, applied to truth and model & \texttt{toy} \\
\texttt{effects.\allowbreak antenna\allowbreak PCO.\allowbreak calibration\allowbreak Residual.\allowbreak enable} & Gate for truth-only phase centre offset mis-calibration & \texttt{true} \\
\texttt{effects.\allowbreak antenna\allowbreak PCO.\allowbreak calibration\allowbreak Residual.\allowbreak receiver\allowbreak Offset\_\allowbreak body\_\allowbreak m} & Receive phase centre mis-calibration per body axis, metres & \texttt{[0.\allowbreak 002,\allowbreak 0.\allowbreak 002,\allowbreak 0.\allowbreak 002]} \\
\texttt{effects.\allowbreak ionosphere.\allowbreak mapping\allowbreak Model} & Obliquity mapping converting vertical to slant ionospheric delay & \texttt{thin\allowbreak Shell} \\
\texttt{effects.\allowbreak ionosphere.\allowbreak shell\allowbreak Height\_\allowbreak m} & Ionospheric thin-shell pierce point height, metres & \texttt{350000} \\
\texttt{effects.\allowbreak solid\allowbreak Earth\allowbreak Tide.\allowbreak truth.\allowbreak enable} & Displaces true tower positions by the solid Earth tide & \texttt{false} \\
\texttt{effects.\allowbreak tower\allowbreak Survey.\allowbreak enable} & Master gate for the tower coordinate survey error & \texttt{true} \\
\texttt{effects.\allowbreak tower\allowbreak Survey.\allowbreak model.\allowbreak enable} & Applies the survey error to the estimator's tower positions & \texttt{false} \\
\texttt{effects.\allowbreak tower\allowbreak Survey.\allowbreak sigma\allowbreak ENU\_\allowbreak m} & Tower survey error one-sigma, east north up, metres & \texttt{[0.\allowbreak 01,\allowbreak 0.\allowbreak 01,\allowbreak 0.\allowbreak 03]} \\
\texttt{effects.\allowbreak tower\allowbreak Survey.\allowbreak truth.\allowbreak enable} & Applies the survey error to the true tower positions & \texttt{true} \\
\texttt{frames.\allowbreak eop\allowbreak Model.\allowbreak enable} & Applies published IERS orientation correction inside the estimator & \texttt{true} \\
\texttt{frames.\allowbreak eop\allowbreak Model.\allowbreak polar\allowbreak Motion\_\allowbreak xp\_\allowbreak arcsec} & Predicted pole x used by the estimator, arcseconds & \texttt{0.\allowbreak 1497} \\
\texttt{frames.\allowbreak eop\allowbreak Model.\allowbreak polar\allowbreak Motion\_\allowbreak yp\_\allowbreak arcsec} & Predicted pole y used by the estimator, arcseconds & \texttt{0.\allowbreak 3497} \\
\texttt{frames.\allowbreak eop\allowbreak Model.\allowbreak ut1\allowbreak Rate\_\allowbreak error\_\allowbreak ms\allowbreak Per\allowbreak Day} & Predicted length-of-day rate correction, milliseconds per day & \texttt{0.\allowbreak 04} \\
\texttt{frames.\allowbreak truth\allowbreak Eop.\allowbreak enable} & Applies truth-side Earth orientation error to tower positions & \texttt{true} \\
\texttt{frames.\allowbreak truth\allowbreak Eop.\allowbreak polar\allowbreak Motion\_\allowbreak xp\_\allowbreak arcsec} & True pole x displacement, arcseconds & \texttt{0.\allowbreak 15} \\
\texttt{frames.\allowbreak truth\allowbreak Eop.\allowbreak polar\allowbreak Motion\_\allowbreak yp\_\allowbreak arcsec} & True pole y displacement, arcseconds & \texttt{0.\allowbreak 35} \\
\texttt{frames.\allowbreak truth\allowbreak Eop.\allowbreak ut1\allowbreak Rate\_\allowbreak error\_\allowbreak ms\allowbreak Per\allowbreak Day} & True length-of-day rate error, milliseconds per day & \texttt{0.\allowbreak 05} \\
\texttt{physics.\allowbreak relativity.\allowbreak clock.\allowbreak enable} & Master gate for the relativistic clock rate offset & \texttt{true} \\
\texttt{physics.\allowbreak relativity.\allowbreak clock.\allowbreak model.\allowbreak enable} & Estimator models the relativistic clock rate offset & \texttt{true} \\
\texttt{physics.\allowbreak relativity.\allowbreak clock.\allowbreak truth.\allowbreak enable} & Puts the relativistic clock rate offset in truth & \texttt{true} \\
\hline
\multicolumn{3}{l}{\textbf{Clocks}} \\ \hline
\texttt{clock.\allowbreak gauge.\allowbreak mode} & How the clock datum ambiguity is resolved & \texttt{external\allowbreak Tower\allowbreak Corrections} \\
\texttt{clock.\allowbreak mode} & Which clocks enter the EKF state vector & \texttt{spacecraft\allowbreak Receiver\allowbreak Clock\allowbreak Only} \\
\texttt{clocks.\allowbreak tower.\allowbreak product.\allowbreak add\allowbreak To\allowbreak R} & Charges the product uncertainty to measurement noise covariance & \texttt{true} \\
\texttt{clocks.\allowbreak tower.\allowbreak product.\allowbreak latency\_\allowbreak s} & Delivery latency of the clock correction, seconds & \texttt{5} \\
\texttt{clocks.\allowbreak tower.\allowbreak product.\allowbreak sigma\allowbreak Bias\_\allowbreak m} & Broadcast tower clock bias residual sigma, metres & \texttt{0.\allowbreak 1} \\
\texttt{clocks.\allowbreak tower.\allowbreak product.\allowbreak sigma\allowbreak Drift\_\allowbreak mps} & Broadcast clock drift residual sigma, metres per second & \texttt{0.\allowbreak 001} \\
\texttt{clocks.\allowbreak tower.\allowbreak product.\allowbreak update\allowbreak Interval\_\allowbreak s} & Rebroadcast interval of the clock correction product, seconds & \texttt{30} \\
\texttt{clocks.\allowbreak tower.\allowbreak product.\allowbreak validity\_\allowbreak s} & Validity span of one clock correction, seconds & \texttt{120} \\
\hline
\multicolumn{3}{l}{\textbf{Truth-side error sources}} \\ \hline
\texttt{biases.\allowbreak inter\allowbreak Frequency.\allowbreak code.\allowbreak model.\allowbreak L1\_\allowbreak m} & L1 code bias the estimator corrects, metres & \texttt{0.\allowbreak 3} \\
\texttt{biases.\allowbreak inter\allowbreak Frequency.\allowbreak code.\allowbreak model.\allowbreak L2\_\allowbreak m} & L2 code bias the estimator corrects, metres & \texttt{0.\allowbreak 405} \\
\texttt{biases.\allowbreak inter\allowbreak Frequency.\allowbreak code.\allowbreak truth.\allowbreak L1\_\allowbreak m} & True L1 code group delay bias, metres & \texttt{0.\allowbreak 3} \\
\texttt{biases.\allowbreak inter\allowbreak Frequency.\allowbreak code.\allowbreak truth.\allowbreak L2\_\allowbreak m} & True L2 code group delay bias, metres & \texttt{0.\allowbreak 45} \\
\texttt{errors.\allowbreak hardware\allowbreak Delay.\allowbreak enable} & Master switch for the tower hardware group delay & \texttt{true} \\
\texttt{errors.\allowbreak hardware\allowbreak Delay.\allowbreak model.\allowbreak enable} & Applies a hardware delay correction on the estimator side & \texttt{false} \\
\texttt{errors.\allowbreak hardware\allowbreak Delay.\allowbreak residual\allowbreak Stochastic.\allowbreak enable} & Draws the white hardware delay residual and charges R & \texttt{true} \\
\texttt{errors.\allowbreak hardware\allowbreak Delay.\allowbreak sigma\_\allowbreak m} & Hardware delay calibration residual one sigma in metres & \texttt{0.\allowbreak 05} \\
\texttt{errors.\allowbreak hardware\allowbreak Delay.\allowbreak truth.\allowbreak enable} & Adds the hardware delay to the truth range & \texttt{true} \\
\texttt{errors.\allowbreak inter\allowbreak Antenna\allowbreak Carrier\allowbreak Bias.\allowbreak drift.\allowbreak enable} & Adds thermal drift to the inter-antenna carrier bias & \texttt{false} \\
\texttt{errors.\allowbreak inter\allowbreak Antenna\allowbreak Carrier\allowbreak Bias.\allowbreak enable} & Injects an unknown per-antenna carrier phase line bias & \texttt{true} \\
\texttt{errors.\allowbreak inter\allowbreak Antenna\allowbreak Carrier\allowbreak Bias.\allowbreak per\allowbreak Signal} & Draws an independent bias for each carrier signal & \texttt{true} \\
\texttt{errors.\allowbreak inter\allowbreak Antenna\allowbreak Carrier\allowbreak Bias.\allowbreak sigma\_\allowbreak cycles} & Inter-antenna carrier phase bias one sigma in cycles & \texttt{0.\allowbreak 25} \\
\texttt{errors.\allowbreak ionosphere.\allowbreak enable} & Master switch for the ionospheric delay error source & \texttt{true} \\
\texttt{errors.\allowbreak ionosphere.\allowbreak higher\allowbreak Order.\allowbreak enable} & Add second and third order ionosphere residuals, truth only & \texttt{true} \\
\texttt{errors.\allowbreak ionosphere.\allowbreak model.\allowbreak correction} & Which broadcast ionosphere correction the estimator applies & \texttt{klobuchar} \\
\texttt{errors.\allowbreak ionosphere.\allowbreak model.\allowbreak enable} & Apply the estimator side ionospheric correction & \texttt{true} \\
\texttt{errors.\allowbreak ionosphere.\allowbreak model\allowbreak Type} & Ionosphere formulation, unlocking the stochastic TEC branch & \texttt{tec\allowbreak Gauss\allowbreak Markov} \\
\texttt{errors.\allowbreak ionosphere.\allowbreak scintillation.\allowbreak S4zen} & Zenith amplitude scintillation index, dimensionless & \texttt{0.\allowbreak 3} \\
\texttt{errors.\allowbreak ionosphere.\allowbreak scintillation.\allowbreak enable} & Add ionospheric amplitude scintillation to truth and R & \texttt{true} \\
\texttt{errors.\allowbreak ionosphere.\allowbreak scintillation.\allowbreak model} & Scintillation formulation supplying the fading factor & \texttt{conker} \\
\texttt{errors.\allowbreak ionosphere.\allowbreak scintillation.\allowbreak phase\allowbreak Scint.\allowbreak enable} & Add phase scintillation to the truth carrier & \texttt{true} \\
\texttt{errors.\allowbreak ionosphere.\allowbreak scintillation.\allowbreak phase\allowbreak Scint.\allowbreak sigma\allowbreak Phi\_\allowbreak rad} & Phase scintillation standard deviation, radians & \texttt{0.\allowbreak 2} \\
\texttt{errors.\allowbreak ionosphere.\allowbreak scintillation.\allowbreak phase\allowbreak Scint.\allowbreak tau\_\allowbreak s} & Correlation time of the phase scintillation, seconds & \texttt{1.\allowbreak 5} \\
\texttt{errors.\allowbreak ionosphere.\allowbreak scintillation.\allowbreak tau\_\allowbreak s} & Correlation time of the scintillation amplitude, seconds & \texttt{30} \\
\texttt{errors.\allowbreak ionosphere.\allowbreak sigma\_\allowbreak m} & Declared vertical ionosphere uncertainty charged in R, metres & \texttt{1.\allowbreak 0} \\
\texttt{errors.\allowbreak ionosphere.\allowbreak stochastic.\allowbreak enable} & Add a time correlated residual TEC component & \texttt{true} \\
\texttt{errors.\allowbreak ionosphere.\allowbreak stochastic.\allowbreak sigma\allowbreak VDelay\allowbreak L1\_\allowbreak ss\_\allowbreak m} & Steady state vertical L1 residual delay, metres & \texttt{0.\allowbreak 3} \\
\texttt{errors.\allowbreak ionosphere.\allowbreak stochastic.\allowbreak tau\_\allowbreak s} & Correlation time of the residual TEC process, seconds & \texttt{600} \\
\texttt{errors.\allowbreak ionosphere.\allowbreak topside\allowbreak Fraction} & Fraction of the vertical TEC column the spacecraft sees & \texttt{1} \\
\texttt{errors.\allowbreak ionosphere.\allowbreak truth.\allowbreak diurnal.\allowbreak enable} & Vary the truth vertical TEC over the daily cycle & \texttt{true} \\
\texttt{errors.\allowbreak ionosphere.\allowbreak truth.\allowbreak diurnal.\allowbreak peak\allowbreak Local\allowbreak Time\_\allowbreak h} & Local solar time of the TEC peak, hours & \texttt{14} \\
\texttt{errors.\allowbreak ionosphere.\allowbreak truth.\allowbreak diurnal.\allowbreak vtec\allowbreak Day\_\allowbreak TECU} & Daytime peak vertical TEC, TECU & \texttt{30} \\
\texttt{errors.\allowbreak ionosphere.\allowbreak truth.\allowbreak diurnal.\allowbreak vtec\allowbreak Night\_\allowbreak TECU} & Night-time floor vertical TEC, TECU & \texttt{6} \\
\texttt{errors.\allowbreak ionosphere.\allowbreak truth.\allowbreak enable} & Inject ionospheric delay into the simulated truth observations & \texttt{true} \\
\texttt{errors.\allowbreak multipath.\allowbreak colored\allowbreak GM.\allowbreak elevation\allowbreak Exponent} & Exponent of the inverse-sine elevation multipath envelope & \texttt{1} \\
\texttt{errors.\allowbreak multipath.\allowbreak colored\allowbreak GM.\allowbreak enable} & Selects the coloured Gauss-Markov multipath process over the white model & \texttt{true} \\
\texttt{errors.\allowbreak multipath.\allowbreak colored\allowbreak GM.\allowbreak shared\allowbreak Across\allowbreak Antennas.\allowbreak enable} & Shares one multipath state per tower across all antennas & \texttt{true} \\
\texttt{errors.\allowbreak multipath.\allowbreak colored\allowbreak GM.\allowbreak sigma\allowbreak Code\allowbreak L1\_\allowbreak ss\_\allowbreak m} & Steady-state L1 code multipath one sigma in metres & \texttt{0.\allowbreak 3} \\
\texttt{errors.\allowbreak multipath.\allowbreak colored\allowbreak GM.\allowbreak tau\_\allowbreak s} & Multipath correlation time in seconds & \texttt{60} \\
\texttt{errors.\allowbreak multipath.\allowbreak enable} & Master switch for the code multipath error & \texttt{true} \\
\texttt{errors.\allowbreak multipath.\allowbreak model.\allowbreak enable} & Applies a multipath correction on the estimator side & \texttt{false} \\
\texttt{errors.\allowbreak multipath.\allowbreak truth.\allowbreak enable} & Adds multipath to the truth pseudorange & \texttt{true} \\
\texttt{errors.\allowbreak troposphere.\allowbreak day\allowbreak Of\allowbreak Year} & Day of year driving the Niell mapping seasonal term & \texttt{180} \\
\texttt{errors.\allowbreak troposphere.\allowbreak enable} & Master switch for the tropospheric delay error source & \texttt{true} \\
\texttt{errors.\allowbreak troposphere.\allowbreak model.\allowbreak enable} & Apply the estimator side tropospheric delay correction & \texttt{true} \\
\texttt{errors.\allowbreak troposphere.\allowbreak model.\allowbreak mapping\allowbreak Type} & Elevation mapping function used by the estimator model & \texttt{niell} \\
\texttt{errors.\allowbreak troposphere.\allowbreak model\allowbreak Type} & Troposphere formulation, surface weather hydrostatic plus stochastic wet delay & \texttt{local\allowbreak Weather\allowbreak GM} \\
\texttt{errors.\allowbreak troposphere.\allowbreak sigma\_\allowbreak m} & Declared zenith troposphere uncertainty charged in R, metres & \texttt{0.\allowbreak 15} \\
\texttt{errors.\allowbreak troposphere.\allowbreak stochastic.\allowbreak enable} & Add a time correlated residual zenith wet delay & \texttt{true} \\
\texttt{errors.\allowbreak troposphere.\allowbreak stochastic.\allowbreak sigma\allowbreak Wet\_\allowbreak ss\_\allowbreak m} & Steady state zenith wet delay variability, metres & \texttt{0.\allowbreak 04} \\
\texttt{errors.\allowbreak troposphere.\allowbreak stochastic.\allowbreak tau\_\allowbreak s} & Correlation time of the wet delay process, seconds & \texttt{10800} \\
\texttt{errors.\allowbreak troposphere.\allowbreak truth.\allowbreak enable} & Inject tropospheric delay into the simulated truth observations & \texttt{true} \\
\texttt{errors.\allowbreak troposphere.\allowbreak truth.\allowbreak mapping\allowbreak Type} & Elevation mapping function used on the truth side & \texttt{niell} \\
\hline
\multicolumn{3}{l}{\textbf{Measurements}} \\ \hline
\texttt{carrier\allowbreak Slip.\allowbreak baseline\allowbreak Differenced\allowbreak Mode.\allowbreak enable} & Whether slip metrics are differenced against a reference antenna & \texttt{true} \\
\texttt{carrier\allowbreak Slip.\allowbreak baseline\allowbreak Differenced\allowbreak Mode.\allowbreak reference\allowbreak Antenna} & Index of the antenna used as differencing reference & \texttt{1} \\
\texttt{carrier\allowbreak Slip.\allowbreak common\allowbreak Mode\allowbreak Compensation.\allowbreak enable} & Whether the common-mode jump is removed before slip detection & \texttt{true} \\
\texttt{carrier\allowbreak Slip.\allowbreak common\allowbreak Mode\allowbreak Compensation.\allowbreak min\allowbreak Rows} & Minimum valid carrier rows needed to estimate the common jump & \texttt{4} \\
\texttt{measurement.\allowbreak sigma\allowbreak Floor\_\allowbreak m} & Lower bound applied to any measurement sigma, metres & \texttt{0.\allowbreak 01} \\
\texttt{measurements.\allowbreak carrier.\allowbreak sigma\_\allowbreak m} & Carrier phase measurement noise sigma, metres & \texttt{0.\allowbreak 01} \\
\texttt{measurements.\allowbreak carrier\allowbreak Mode} & Carrier phase mode, off, diagnostic or float ambiguity & \texttt{ekf\allowbreak Float} \\
\texttt{measurements.\allowbreak carrier\allowbreak Phase.\allowbreak enable} & Legacy carrier phase gate, a fallback when carrier mode absent & \texttt{true} \\
\texttt{measurements.\allowbreak code\allowbreak Noise.\allowbreak cn0.\allowbreak base\_\allowbreak d\allowbreak BHz} & Baseline carrier-to-noise density in the C/N0 link model, dB-Hz & \texttt{45} \\
\texttt{measurements.\allowbreak code\allowbreak Noise.\allowbreak cn0.\allowbreak elevation\allowbreak Gain\_\allowbreak d\allowbreak B} & Carrier-to-noise density gain from horizon to zenith, decibels & \texttt{6} \\
\texttt{measurements.\allowbreak code\allowbreak Noise.\allowbreak cn0.\allowbreak sigma\allowbreak At45d\allowbreak BHz\_\allowbreak m} & Reference code noise sigma at 45 dB-Hz, metres & \texttt{0.\allowbreak 3} \\
\texttt{measurements.\allowbreak code\allowbreak Noise.\allowbreak model} & Code noise model, constant, elevation weighted or C/N0 & \texttt{cn0} \\
\texttt{measurements.\allowbreak doppler.\allowbreak enable} & Whether Doppler range-rate observables are generated & \texttt{true} \\
\texttt{measurements.\allowbreak doppler.\allowbreak sigma\_\allowbreak mps} & Doppler measurement noise sigma, metres per second & \texttt{0.\allowbreak 0424} \\
\texttt{measurements.\allowbreak doppler.\allowbreak use\allowbreak In\allowbreak EKF} & Whether Doppler rows update the filter & \texttt{true} \\
\texttt{measurements.\allowbreak two\allowbreak Way\allowbreak Time\allowbreak Transfer.\allowbreak conservative\allowbreak Product\allowbreak Correlation} & Whether tower clock product variance is inflated for within-interval correlation & \texttt{true} \\
\texttt{measurements.\allowbreak two\allowbreak Way\allowbreak Time\allowbreak Transfer.\allowbreak enable} & Whether tower-spacecraft two-way time transfer rows are generated & \texttt{true} \\
\texttt{measurements.\allowbreak two\allowbreak Way\allowbreak Time\allowbreak Transfer.\allowbreak include\allowbreak Reciprocity\allowbreak Residual} & Whether motion non-reciprocity of the two-way path is modelled & \texttt{true} \\
\texttt{measurements.\allowbreak two\allowbreak Way\allowbreak Time\allowbreak Transfer.\allowbreak mode} & Two-way observable formulation, first-order reciprocal or four-timestamp & \texttt{first\allowbreak Order\allowbreak Reciprocal} \\
\texttt{measurements.\allowbreak two\allowbreak Way\allowbreak Time\allowbreak Transfer.\allowbreak reciprocity\allowbreak Sigma\_\allowbreak m} & Residual non-reciprocity sigma of the two-way path, metres & \texttt{0.\allowbreak 005} \\
\texttt{measurements.\allowbreak two\allowbreak Way\allowbreak Time\allowbreak Transfer.\allowbreak sigma\_\allowbreak m} & Two-way time transfer noise sigma, metres & \texttt{0.\allowbreak 03} \\
\texttt{measurements.\allowbreak two\allowbreak Way\allowbreak Time\allowbreak Transfer.\allowbreak towers} & Which ground towers are two-way capable & \texttt{all} \\
\texttt{measurements.\allowbreak two\allowbreak Way\allowbreak Time\allowbreak Transfer.\allowbreak use\allowbreak In\allowbreak EKF} & Whether two-way time transfer rows update the filter & \texttt{true} \\
\texttt{measurements.\allowbreak two\allowbreak Way\allowbreak Time\allowbreak Transfer.\allowbreak warmup\_\allowbreak s} & Delay before two-way rows begin, seconds & \texttt{0} \\
\hline
\multicolumn{3}{l}{\textbf{Estimator configuration}} \\ \hline
\texttt{estimation.\allowbreak ambiguity.\allowbreak initial\allowbreak Sigma\_\allowbreak m} & Initial one-sigma prior on each carrier ambiguity state, metres & \texttt{100} \\
\texttt{estimation.\allowbreak ambiguity\allowbreak Mode} & How carrier phase ambiguities are parameterised as filter states & \texttt{float\allowbreak Per\allowbreak Tower\allowbreak Receiver\allowbreak Signal} \\
\texttt{estimation.\allowbreak ionosphere\allowbreak Mode} & Whether and how the filter estimates ionospheric delay states & \texttt{per\allowbreak Tower\allowbreak Slant} \\
\texttt{estimation.\allowbreak slant\allowbreak Iono.\allowbreak initial\allowbreak Sigma\_\allowbreak m} & Initial one-sigma prior on the slant-ionosphere state, metres & \texttt{5.\allowbreak 0} \\
\texttt{estimation.\allowbreak slant\allowbreak Iono.\allowbreak sigma\_\allowbreak ss\_\allowbreak m} & Steady-state sigma of the slant-ionosphere state, metres & \texttt{1.\allowbreak 0} \\
\texttt{estimation.\allowbreak slant\allowbreak Iono.\allowbreak tau\_\allowbreak s} & Correlation time of the slant-ionosphere Gauss-Markov state, seconds & \texttt{600} \\
\texttt{estimation.\allowbreak tropo\allowbreak Zwd.\allowbreak initial\allowbreak Sigma\_\allowbreak m} & Initial one-sigma prior on the zenith wet delay state, metres & \texttt{0.\allowbreak 1} \\
\texttt{estimation.\allowbreak tropo\allowbreak Zwd.\allowbreak sigma\_\allowbreak ss\_\allowbreak m} & Steady-state sigma of the zenith wet delay state, metres & \texttt{0.\allowbreak 04} \\
\texttt{estimation.\allowbreak tropo\allowbreak Zwd.\allowbreak tau\_\allowbreak s} & Correlation time of the zenith wet delay state, seconds & \texttt{10800} \\
\texttt{estimation.\allowbreak troposphere\allowbreak Mode} & Whether and how the filter estimates wet tropospheric delay states & \texttt{per\allowbreak Tower\allowbreak Zwd} \\
\texttt{estimator.\allowbreak attitude.\allowbreak use\allowbreak Carrier\allowbreak Partials} & Lets carrier phase rows contribute attitude partials to the filter & \texttt{false} \\
\texttt{estimator.\allowbreak attitude.\allowbreak use\allowbreak Code\allowbreak Partials} & Lets pseudorange rows contribute attitude partials to the filter & \texttt{false} \\
\texttt{estimator.\allowbreak attitude\allowbreak Carrier\allowbreak Mode} & Selects the carrier phase attitude measurement path & \texttt{off} \\
\texttt{estimator.\allowbreak attitude\allowbreak Init\allowbreak Mode} & Selects the absolute attitude seeding search before tracking & \texttt{none} \\
\texttt{estimator.\allowbreak dynamics.\allowbreak perturbations.\allowbreak epoch\allowbreak JD\_\allowbreak TT} & Julian date in terrestrial time anchoring the Sun Moon ephemeris & \texttt{2451545.\allowbreak 0} \\
\texttt{estimator.\allowbreak dynamics.\allowbreak perturbations.\allowbreak luni\allowbreak Solar.\allowbreak enable} & Sun and Moon third body acceleration in the filter propagator & \texttt{true} \\
\texttt{estimator.\allowbreak dynamics.\allowbreak perturbations.\allowbreak srp.\allowbreak Cr} & Filter reflectivity coefficient for solar radiation pressure, dimensionless & \texttt{1.\allowbreak 24} \\
\texttt{estimator.\allowbreak dynamics.\allowbreak perturbations.\allowbreak srp.\allowbreak area\allowbreak To\allowbreak Mass\_\allowbreak m2pkg} & Filter area to mass ratio in square metres per kilogram & \texttt{0.\allowbreak 02} \\
\texttt{estimator.\allowbreak dynamics.\allowbreak perturbations.\allowbreak srp.\allowbreak enable} & Solar radiation pressure in the filter propagator & \texttt{true} \\
\texttt{estimator.\allowbreak elevation\allowbreak Mask\_\allowbreak rad} & Elevation cutoff in radians below which measurements are rejected & \texttt{0.\allowbreak 1745\allowbreak 3292\allowbreak 5199\allowbreak 43295} \\
\texttt{estimator.\allowbreak empirical\allowbreak Accel.\allowbreak enable} & Master gate for three empirical radial along track cross track accelerations & \texttt{false} \\
\texttt{estimator.\allowbreak empirical\allowbreak Accel.\allowbreak use\allowbreak In\allowbreak EKF} & Appends the empirical acceleration states to the filter & \texttt{false} \\
\texttt{estimator.\allowbreak imu.\allowbreak filter.\allowbreak P0\_\allowbreak bias\_\allowbreak radps} & Initial gyro bias uncertainty one sigma in radians per second & \texttt{3e-05} \\
\texttt{estimator.\allowbreak imu.\allowbreak filter.\allowbreak arw\_\allowbreak rad\_\allowbreak per\_\allowbreak sqrt\_\allowbreak s} & Filter gyro angle random walk in radians per root second & \texttt{0.\allowbreak 0002} \\
\texttt{estimator.\allowbreak imu.\allowbreak filter.\allowbreak rrw\_\allowbreak rad\_\allowbreak per\_\allowbreak s\_\allowbreak sqrt\_\allowbreak s} & Filter gyro rate random walk, radians per second per root second & \texttt{3e-06} \\
\texttt{estimator.\allowbreak imu.\allowbreak truth.\allowbreak arw\_\allowbreak rad\_\allowbreak per\_\allowbreak sqrt\_\allowbreak s} & Truth gyro angle random walk in radians per root second & \texttt{0.\allowbreak 0002} \\
\texttt{estimator.\allowbreak imu.\allowbreak truth.\allowbreak bias0\allowbreak Sigma\_\allowbreak radps} & Truth gyro initial bias draw one sigma in radians per second & \texttt{3e-05} \\
\texttt{estimator.\allowbreak imu.\allowbreak truth.\allowbreak rrw\_\allowbreak rad\_\allowbreak per\_\allowbreak s\_\allowbreak sqrt\_\allowbreak s} & Truth gyro rate random walk, radians per second per root second & \texttt{3e-06} \\
\texttt{estimator.\allowbreak lambda.\allowbreak enable} & Master gate for LAMBDA integer ambiguity resolution & \texttt{false} \\
\texttt{estimator.\allowbreak process\allowbreak Noise.\allowbreak model\allowbreak Mismatch.\allowbreak enable} & Adds extra process noise for a truth versus filter force gap & \texttt{false} \\
\texttt{estimator.\allowbreak sigma\_\allowbreak accel\_\allowbreak mps2} & Residual acceleration process noise one sigma in metres per second squared & \texttt{1e-06} \\
\texttt{estimator.\allowbreak srp\allowbreak Coefficient.\allowbreak enable} & Master gate for the solar radiation pressure scale state & \texttt{false} \\
\texttt{estimator.\allowbreak srp\allowbreak Coefficient.\allowbreak use\allowbreak In\allowbreak EKF} & Appends the solar pressure scale state to the filter & \texttt{false} \\
\texttt{estimator.\allowbreak star\allowbreak Tracker.\allowbreak white\allowbreak Angular\allowbreak Sigma\_\allowbreak rad} & Star tracker white angular noise one sigma in radians & \texttt{0.\allowbreak 0001\allowbreak 4544\allowbreak 4104\allowbreak 3328608} \\
\hline
\multicolumn{3}{l}{\textbf{Ensemble reporting}} \\ \hline
\texttt{report.\allowbreak monte\allowbreak Carlo.\allowbreak confidence} & Confidence level of the chi-squared consistency band & \texttt{0.\allowbreak 99} \\
\texttt{report.\allowbreak monte\allowbreak Carlo.\allowbreak duration\_\allowbreak s} & Arc length of each Monte Carlo run, seconds & \texttt{900} \\
\texttt{report.\allowbreak monte\allowbreak Carlo.\allowbreak enable} & Appends a pooled multi-seed NEES and NIS consistency check & \texttt{true} \\
\texttt{report.\allowbreak monte\allowbreak Carlo.\allowbreak n\allowbreak Seeds} & Number of seeded runs in the Monte Carlo ensemble & \texttt{12} \\
\hline
\end{longtable}
\end{center}


\end{document}
